% !TEX program = latexmk
\documentclass[paper=8.125in:10.250in,pagesize=pdftex,
    headinclude=false,footinclude=true,final,oneside,sfdefaults=false]{kaobook}

\def\sagdir{sag}
\def\rootdir{../../..}

\input{\rootdir/utils/preamble}
%\input{\rootdir/utils/body-init}

%\title{ACT4E Exam \\
%\normalsize{ETH Zurich, Summer 2021}}

%\date{}


\renewcommand{\marginlayout}{%
    \newgeometry{
        top=22.4mm,             % height of the top margin
        bottom=22.4mm,          % height of the bottom margin
        inner=20.8mm,    % width of the inner margin
        outer=20.8mm,            % width of the outer margin
        %textwidth=170mm,        % width of the text
        %marginparsep=8.2mm,     % width between text and margin
        marginparwidth=05.4mm,  % width of the margin
    }%
}



\begin{document}
    %\setboolean{bookversion}{true}
    \setboolean{statuscolors}{false}
    \setboolean{debugimages}{false}
    \setboolean{cachepdf}{false}
    \setboolean{devel}{false}
    \setboolean{showslides}{false}
    \setboolean{instructors}{true}
    \setboolean{codeexercises}{true}

    \pagestyle{myheadings}
    \markright{Applied Category Theory for Engineering I \hfill ETH Zurich, Fall 2026}

    \setcounter{chapter}{2}

    \begin{center}
        {\large{\textbf{Homework sheet \thechapter}}} \\
    \end{center}

    \begin{homework}[\exname{DivisibilityOf72}] [3 points]
Consider the set $\setA$ of natural numbers which divide the natural number 72,
\begin{equation}
\setA = \makeset{ n \setin \natnumbers \mid \exists k \setin \natnumbers \colon nk = 72},
\end{equation}
equipped with the partial order ``$\posleq$'' defined by
\begin{equation}
\ela \posleq \elb \quad \text{ if, and only if } \quad \ela \text{ divides } \elb
\end{equation}
for all $\ela, \elb \setin \setA$. Your task: draw a Hasse diagram of this partially ordered set.
\end{homework}
\solutionof{DivisibilityOf72}


%    \begin{homework}[\exname{PolynomialDivisibility}] [5 points]
%        \label{ex:PolynomialDivisibility}
%
%        Let~$\setA$ be the set of all polynomials with coefficients in $\reals$:
%        \begin{equation}
%\setA = \makeset{ a_nX^n + a_{n-1}X^{n-1} + \dots + a_1X + a_0 \mid n \setin \natnumbers, 0 \leq i \leq n, a_i \setin \reals }.
%\end{equation}
%(The symbol $X$ here is a ``formal variable'' -- it is a sort of placeholder).
%
%
%        Recall that a polynomial~$p$ \emph{divides} a polynomial~$q$ if there exists a polynomial~$m$ such that~$p \cdot m = q$.
%        If~$p$ divides~$q$ we denote this by~$p \vert q$.
%
%        Divisibility defines an endorelation $\relA$ on~$\setA$ by saying~$p$ is related to~$q$ iff~$p \vert q$:
%        \begin{equation}
%        \prfdoubleperiod{
%p \relA q
%}{
%p \vert q
%}
%\end{equation}
%
%        Does $\relA$ define a pre-order structure on~$\setA$?
%        Does $\relA$ define a poset structure on~$\setA$?
%        Justify your answers.
%    \end{homework}
%    \solutionof{PolynomialDivisibility}



% \begin{homework}[\exname{ProductPoset}] [4 points]
%
%Consider the posets~$\posA = \tup{\posAset, \posleq_\posA}$ and~$\posB = \tup{\posBset, \posleq_\posB}$ described respectively by the following Hasse diagrams.
%
%        \begin{equation}
%            %\includegraphics[width=0.5\linewidth]{pics/ExamPosets.png}
%            \includesag{exam_posets}
%        \end{equation}
%
%In this exercise, we are interested in the product poset $\posA \Ptimes \posB$. You do not need to prove your answers to the following questions (simply state your answer).
%
%        \begin{enumerate}
%        \item How many elements does the underlying set of $\posA \Ptimes \posB$ have?
%        \item Is $\posA \Ptimes \posB$ a total order?
%        \item Draw a Hasse diagram of this poset.
%        \end{enumerate}
%
%    \end{homework}
%\solutionof{ProductPoset}


\begin{homework}[\exname{WhichMapsMonotone}] [4 points]

Consider the posets~$\posA = \tup{\posAset, \posleq_\posA}$ and~$\posB = \tup{\posBset, \posleq_\posB}$ described respectively by the following Hasse diagrams.

        \begin{center}
            %\includegraphics[width=0.5\linewidth]{pics/ExamPosets.png}
            \includesag{exam_posets}
        \end{center}

        The following diagrams show functions~$\posAset \sto \posBset$.
        We will call them~$\mapa_1$,~$\mapa_2$,~$\mapa_3$ and~$\mapa_4$, respectively.

        \begin{center}
            \setlength{\tabcolsep}{8pt}
            \begin{tabular}{cc@{\hspace{40pt}}cc}
                $\mapa_1$ &
                %\includegraphics[width=0.5\linewidth]{pics/ExamPosetMap1.png}
                \includesag{exam_f1} &
                $\mapa_2$ &
                %\includegraphics[width=0.5\linewidth]{pics/ExamPosetMap2.png}
                \includesag{exam_f2} \\[+40pt]
                $\mapa_3$ &
                %\includegraphics[width=0.5\linewidth]{pics/ExamPosetMap3.png}
                \includesag{exam_f3} &
                $\mapa_4$ &
                %\includegraphics[width=0.5\linewidth]{pics/ExamPosetMap4.png}
                \includesag{exam_f4}
            \end{tabular}
        \end{center}

Which of the functions~$\mapa_1$,~$\mapa_2$,~$\mapa_3$ and~$\mapa_4$ are monotone? You do not need to prove your answers.
\end{homework}
\solutionof{WhichMapsMonotone}


%    \begin{homework}[\exname{MonotoneMapCheck}] [6 points]
%        \label{ex:MonotoneMapCheck}
%
%        Prove your answers to the following questions.
%        \begin{enumerate}
%            \item Is the function $\mapa \colon \tup{\wnumbers, \leq} \to \tup{\wnumbers, \leq}, \ela \mapsto \ela^2$ monotone?
%            \item Let $\setA = \makeset{a, b, c}$ and consider the posets $\tup{\powerset{\setA}, \subseteq}$ and $\tup{\natnumbers, \leq}$.
%                  Let $\mapa \colon \powerset{\setA} \to \natnumbers, \subA \mapsto \vert \subA \vert$ be the function which computes the cardinality of any subset of $\setA$.
%                  Is $\mapa$ monotone?
%            \item Consider the set $\posAset$ of natural numbers which divide the number 36, equipped with the partial order ``$\preccurlyeq$'' such that $\ela \preccurlyeq \elb$ if and only if $\ela$ divides $\elb$.
%                  Call this poset $\posA = \tup{\posAset, \preccurlyeq}$, and let $\mapa \colon \posAset \to \makeset{\false, \true}$ be defined by
%                  \begin{equation}
%                      \mapa(\ela) =
%                      \begin{cases}
%                          \true  & \text{ if } \ela \text{ is an even number, } \\
%                          \false & \text{ if } \ela \text{ is an odd number.
%                          }
%                      \end{cases}
%                  \end{equation}
%                  Is $\mapa$ monotone if we equip $\makeset{\false, \true}$ with the partial order such that $\false \leq \true$?
%        \end{enumerate}
%    \end{homework}
%    \solutionof{MonotoneMapCheck}

    \begin{homework}[\exname{UpperSetsAsMonotoneMaps}][6 points]
    \label{ex:UpperSetsAsMonotoneMaps}

    In this exercise, $\Bool$ denotes the poset with two elements,~$\true$ and~$\false$, whose only non-trivial inequality is~$\false \leq \true$.

    Let~$\posA = \tup{\posAset, \posAleq}$ be a poset.
    \begin{enumerate}
        \item Prove that for any monotone map~$\mapa \colon \posA \mto \Bool$, the subset
        \begin{equation}
            \subA_{\mapa} \definedas \makeset{ \posAel \setin \posAset \mid \mapa(\posAel) = \true } \setsubseteq \posAset
        \end{equation}
        is an \SY{upper set}.
        \item Conversely, prove that every upper set~$\subA \setsubseteq \posAset$ is of this form for a \emph{unique} monotone map; that is, there exists exactly one monotone map~$\mapa \colon \posA \mto \Bool$ such that~$\subA_{\mapa} = \subA$.
    \end{enumerate}
\end{homework}
\solutionof{UpperSetsAsMonotoneMaps}




\begin{homework}[\exname{CurryingDesignProblemsBasic}][5 points]
    \label{ex:CurryingDesignProblemsBasic}

    Let~$\adp\colon \funspop \Ptimes \ressp\toinPos \Bool$ be a design problem, and let~$\uppersets \ressp$ denote the set of upper sets of~$\ressp$.
    \begin{enumerate}
        \item Let~$\fun \setin \funspop$.
        Prove that the map~$\ressp \mto \Bool$, $\res \mapsto \adp(\tup{\fun, \res})$ is monotone, and use \cref{ex:UpperSetsAsMonotoneMaps} to conclude that the subset
        \begin{equation}
            \widehat{\adp}(\fun) \definedas \makeset{ \res \setin \ressp \mid \adp(\tup{\fun, \res}) = \true} \setsubseteq \ressp
        \end{equation}
        is an upper set of~$\ressp$.
        \item By the previous part, we have a function
        \begin{equation}
            \label{eq:dp-left-curry}
            \defmap{ \widehat{\adp}
            }{ \funspop
            }{ \mto
            }{ \tup{\uppersets \ressp, \setsubseteq}
            }{ \fun
            }{ \makeset{ \res \setin \ressp \mid \adp(\tup{\fun, \res}) = \true}
            }
        \end{equation}
        Prove that~$\widehat{\adp} \colon \funspop \mto \tup{\uppersets \ressp, \setsubseteq}$ is a monotone map, where ``$\setsubseteq$'' is the ordering by inclusion of subsets.
        \item Deduce that for all~$\fun, \EM{{\colF g}} \setin \funsp$, if~$\fun \posleq_{\funsp} \EM{{\colF g}}$, then~$\widehat{\adp}(\EM{{\colF g}}) \setsubseteq \widehat{\adp}(\fun)$.
        Explain in one sentence what this means for a designer.
    \end{enumerate}
\end{homework}
\solutionof{CurryingDesignProblemsBasic}






\end{document}
